\chapter{2006年陕西卷（理）}
\section{单选题}
\begin{problem}
已知集合 \(P = \{x \in \N \mid 1 \leqslant x \leqslant 10\}\)，集合 \(Q = \{x \in \R \mid x^2 + x - 6 \leqslant 0\}\)，则 \(P \cap Q\) 等于（\quad）

\choices
    {\(\{2\}\)}
    {\(\{1, 2\}\)}
    {\(\{2,3\}\)}
    {\(\{1,2,3\}\)}
\end{problem}

\begin{answer}
B
\end{answer}

\begin{problem}
复数 \(\frac{(1 + \i)^2}{1 - \i}\) 等于（\quad）

\choices
    {\(1 - \i\)}
    {\(1 + \i\)}
    {\(-1 + \i\)}
    {\(-1 - \i\)}
\end{problem}

\begin{answer}
C
\end{answer}

\begin{problem}
\(\displaystyle\lim_{n\to \infty}\frac{1}{2n(\sqrt{n^2 + 1} - \sqrt{n^2 - 1})}\) 等于（\quad）

\choices
    {\(1\)}
    {\(\frac{1}{2}\)}
    {\(\frac{1}{4}\)}
    {\(0\)}
\end{problem}

\begin{answer}
B
\end{answer}

\begin{problem}
设函数 \( f(x) = \log_a (x + b) \)（\(a > 0\)，\(a \neq 1\)）的图象过点 \((2,1)\)，其反函数的图象过点 \((2,8)\)，则 \( a + b \) 等于（\quad）

\choices
    {\(6\)}
    {\(5\)}
    {\(4\)}
    {\(3\)}
\end{problem}

\begin{answer}
C
\end{answer}

\begin{problem}
设直线过点 \((0, a)\)，其斜率为 1，且与圆 \(x^{2} + y^{2} = 2\) 相切，则 \(a\) 的值为（\quad）

\choices
    {\(\pm \sqrt{2}\)}
    {\(\pm 2\)}
    {\(\pm 2\sqrt{2}\)}
    {\(\pm 4\)}
\end{problem}

\begin{answer}
B
\end{answer}

\begin{problem}
“等式 \( \sin(\alpha+\gamma)=\sin2\beta \) 成立”是“ \(\alpha\)，\(\beta\)，\(\gamma\) 成等差数列”的（\quad）

\choices
    {必要而不充分条件}
    {充分而不必要条件}
    {充分必要条件}
    {既不充分也不必要条件}
\end{problem}

\begin{answer}
A
\end{answer}

\begin{problem}
已知双曲线 \(\frac{x^2}{a^2} - \frac{y^2}{2} = 1\)（\(a > \sqrt{2}\)）的两条渐近线的夹角为 \(\frac{\pi}{3}\)，则双曲线的离心率为（\quad）

\choices
    {\(2\)}
    {\(\sqrt{3}\)}
    {\(\frac{2\sqrt{6}}{3}\)}
    {\(\frac{2\sqrt{3}}{3}\)}
\end{problem}

\begin{answer}
D
\end{answer}

\begin{problem}
已知不等式 \((x + y)\left(\frac{1}{x} + \frac{a}{y}\right) \geqslant 9\) 对任意正实数 \(x\)，\(y\) 恒成立，则正实数 \(a\) 的最小值为（\quad）

\choices
    {\(2\)}
    {\(4\)}
    {\(6\)}
    {\(8\)}
\end{problem}

\begin{answer}
B
\end{answer}

\begin{problem}
已知非零向量 \(\overrightarrow{AB}\) 与 \(\overrightarrow{AC}\) 满足 \(\left(\frac{\overrightarrow{AB}}{|\overrightarrow{AB}|} + \frac{\overrightarrow{AC}}{|\overrightarrow{AC}|}\right)\cdot \overrightarrow{BC} = 0\) 且 \(\frac{\overrightarrow{AB}}{|\overrightarrow{AB}|}\cdot \frac{\overrightarrow{AC}}{|\overrightarrow{AC}|} = \frac{1}{2}\)，则 \(\triangle ABC\) 为（\quad）

\choices
    {三边均不相等的三角形}
    {直角三角形}
    {等腰非等边三角形}
    {等边三角形}
\end{problem}

\begin{answer}
D
\end{answer}

\begin{problem}
已知函数 \( f(x) = ax^2 + 2ax + 4 \)（\(0 < a < 3\)）. 若 \(x_1 < x_2\)，\(x_1 + x_2 = 1 - a\)，则（\quad）

\choices
    {\(f(x_{1}) < f(x_{2})\)}
    {\(f(x_{1}) = f(x_{2})\)}
    {\(f(x_{1}) > f(x_{2})\)}
    {\(f(x_{1})\) 与 \(f(x_{2})\) 的大小不能确定}
\end{problem}

\begin{answer}
A
\end{answer}

\begin{problem}
已知平面 \(\alpha\) 外不共线的三点 \(A\)，\(B\)，\(C\) 到 \(\alpha\) 的距离都相等. 则正确的结论是（\quad）

\choices
    {平面 \(ABC\) 必平行于 \(\alpha\)}
    {平面 \(ABC\) 必与 \(\alpha\) 相交}
    {平面 \(ABC\) 必不垂直于 \(\alpha\)}
    {存在 \(\triangle ABC\) 的一条中位线平行于 \(\alpha\) 或在 \(\alpha\) 内}
\end{problem}

\begin{answer}
D
\end{answer}

\begin{problem}
为确保信息安全，信息需加密传播. 发送方由明文 \(\rightarrow\) 密文（加密），接收方由密文 \(\rightarrow\) 明文（解密）. 已知加密规则为：明文 \(a\)，\(b\)，\(c\)，\(d\) 对应密文 \(a + 2b\)，\(2b + c\)，\(2c + 3d\)，\(4d\). 例如，明文 \(1\)，\(2\)，\(3\)，\(4\) 对应密文 \(5\)，\(7\)，\(18\)，\(16\). 当接收方收到密文 \(14\)，\(9\)，\(23\)，\(28\) 时，则解密得到的明文为（\quad）

\choices
    {\(4,6,1,7\)}
    {\(7,6,1,4\)}
    {\(6,4,1,7\)}
    {\(1,6,4,7\)}
\end{problem}

\begin{answer}
C
\end{answer}

\section{填空题}
\begin{problem}
\(\cos 43^{\circ} \cos 77^{\circ} + \sin 43^{\circ} \cos 167^{\circ}\) 的值为\fillinblank{}.
\end{problem}

\begin{answer}
\(-\frac{1}{2}\)
\end{answer}

\begin{problem}
\(\left(3x - \frac{1}{\sqrt{x}}\right)^{12}\) 展开式中 \(x^{-3}\) 的系数为\fillinblank{}.（用数字作答）
\end{problem}

\begin{answer}
\(594\)
\end{answer}

\begin{problem}
水平桌面 \(\alpha\) 上放有 4 个半径均为 \(2R\) 的球，且相邻的球都相切（球心的连线构成正方形）. 在这 4 个球的上面放 1 个半径为 \(R\) 的小球，它和下面的 4 个球恰好都相切，则小球的球心到水平桌面 \(\alpha\) 的距离是\fillinblank{}.
\end{problem}

\begin{answer}
\(3R\)
\end{answer}

\begin{problem}
某校从 8 名教师中选派 4 名教师同时去 4 个边远地区支教（每地 1 人），其中甲和乙不同去，甲和丙只能同去或同不去，则不同的选派方案共有\fillinblank{}种.
\end{problem}

\begin{answer}
\(600\)
\end{answer}

\section{解答题}
\begin{problem}
已知函数 \( f(x) = \sqrt{3} \sin \left(2x - \frac{\pi}{6}\right) + 2 \sin^2 \left(x - \frac{\pi}{12}\right) \)（\(x \in \R\)）.

\begin{enumerate}
\item 求函数 \( f(x) \) 的最小正周期；
\item 求使函数 \( f(x) \) 取得最大值的 \(x\) 的集合.
\end{enumerate}
\end{problem}

\begin{solution}
\begin{enumerate}
\item 令 \(u=2x-\frac{\pi}{6}\)，则 \(2\sin^2\left(x-\frac{\pi}{12}\right)=1-\cos\left(2x-\frac{\pi}{6}\right)=1-\cos u\)，所以 \(f(x)=1+\sqrt{3}\sin u-\cos u=1+2\sin\left(u-\frac{\pi}{6}\right)=1+2\sin\left(2x-\frac{\pi}{3}\right)\). 因此，函数 \(f(x)\) 的最小正周期为 \(\pi\).
\item 当且仅当 \(\sin\left(2x-\frac{\pi}{3}\right)=1\) 时，\(f(x)\) 取得最大值，即 \(2x-\frac{\pi}{3}=\frac{\pi}{2}+2k\pi\)，其中 \(k\in\Z\). 解得 \(x=\frac{5\pi}{12}+k\pi\)，所以使函数取得最大值的 \(x\) 的集合为 \(\left\{x\mid x=\frac{5\pi}{12}+k\pi,k\in\Z\right\}\).
\end{enumerate}
\end{solution}

\begin{problem}
甲，乙，丙 3 人投篮，投进的概率分别是 \(\frac{1}{3}\)，\(\frac{2}{5}\)，\(\frac{1}{2}\).

\begin{enumerate}
\item 现 3 人各投篮 1 次，求 3 人都没有投进的概率；
\item 用 \( \xi \) 表示乙投篮 3 次的进球数，求随机变量 \( \xi \) 的概率分布及数学期望 \( E(\xi) \).
\end{enumerate}
\end{problem}

\begin{solution}
\begin{enumerate}
\item 甲、乙、丙分别没有投进的概率为 \(1-\frac{1}{3}=\frac{2}{3}\)、\(1-\frac{2}{5}=\frac{3}{5}\)、\(1-\frac{1}{2}=\frac{1}{2}\). 由于 3 人的投篮结果相互独立，3 人都没有投进的概率为 \(\frac{2}{3}\cdot\frac{3}{5}\cdot\frac{1}{2}=\frac{1}{5}\).
\item 乙每次投篮投进的概率为 \(\frac{2}{5}\)，投失的概率为 \(\frac{3}{5}\)，所以 \(\xi\) 服从参数为 \(3\)，\(\frac{2}{5}\) 的二项分布. 对 \(k\in\{0,1,2,3\}\)，有 \(P(\xi=k)=\mathrm C_3^k\left(\frac{2}{5}\right)^k\left(\frac{3}{5}\right)^{3-k}\). 因而 \(P(\xi=0)=\frac{27}{125}\)，\(P(\xi=1)=\frac{54}{125}\)，\(P(\xi=2)=\frac{36}{125}\)，\(P(\xi=3)=\frac{8}{125}\)，这就是随机变量 \(\xi\) 的概率分布. 于是 \(E(\xi)=0\cdot\frac{27}{125}+1\cdot\frac{54}{125}+2\cdot\frac{36}{125}+3\cdot\frac{8}{125}=\frac{150}{125}=\frac{6}{5}\).
\end{enumerate}
\end{solution}

\begin{problem}
如图，\(\alpha \perp \beta\)，\(\alpha \cap \beta = l\)，\(A \in \alpha\)，\(B \in \beta\)，点 \(A\) 在直线 \(l\) 上的射影为 \(A_1\)，点 \(B\) 在 \(l\) 上的射影为 \(B_1\). 已知 \(AB = 2\)，\(AA_1 = 1\)，\(BB_1 = \sqrt{2}\). 求：

\begin{enumerate}
\item 直线 \(AB\) 分别与平面 \(\alpha\)，\(\beta\) 所成角的大小；
\item 二面角 \(A_{1} - AB - B_{1}\) 的大小.
\end{enumerate}

\bitmapfigure[max width=0.38\linewidth]{img/2006/shaanxi_science/q19_fig1.png}
\end{problem}

\begin{solution}
\begin{enumerate}
\item 建立直角坐标系，使直线 \(l\) 为 \(x\) 轴，平面 \(\alpha\) 为 \(xOy\) 平面，平面 \(\beta\) 为 \(xOz\) 平面. 取 \(A_1=(0,0,0)\)，则可设 \(A=(0,1,0)\)；设 \(B_1=(d,0,0)\)，则 \(B=(d,0,\sqrt{2})\). 由 \(AB=2\) 得 \(d^2+1^2+(\sqrt{2})^2=2^2\)，所以 \(d^2=1\). 适当选取坐标轴正方向，可令 \(d=1\). 直线 \(AB\) 垂直于平面 \(\alpha\) 的分量为 \(BB_1\)，垂直于平面 \(\beta\) 的分量为 \(AA_1\)，故设直线 \(AB\) 分别与平面 \(\alpha\)，\(\beta\) 所成的角为 \(\theta\)，\(\varphi\)，则 \(\sin\theta=\frac{BB_1}{AB}=\frac{\sqrt{2}}{2}\)，\(\sin\varphi=\frac{AA_1}{AB}=\frac{1}{2}\). 因而 \(\theta=\frac{\pi}{4}\)，\(\varphi=\frac{\pi}{6}\).
\item 由上述坐标可得 \(\overrightarrow{A_1A}=(0,1,0)\)，\(\overrightarrow{A_1B}=(1,0,\sqrt{2})\)，\(\overrightarrow{B_1A}=(-1,1,0)\)，\(\overrightarrow{B_1B}=(0,0,\sqrt{2})\). 于是，平面 \(A_1AB\) 和平面 \(B_1AB\) 的一个法向量分别为 \(\bs{n}_1=\overrightarrow{A_1A}\times\overrightarrow{A_1B}=(\sqrt{2},0,-1)\)，\(\bs{n}_2=\overrightarrow{B_1A}\times\overrightarrow{B_1B}=(\sqrt{2},\sqrt{2},0)\). 设二面角 \(A_1-AB-B_1\) 为 \(\psi\)，则 \(\cos\psi=\frac{|\bs{n}_1\cdot\bs{n}_2|}{|\bs{n}_1||\bs{n}_2|}=\frac{2}{\sqrt{3}\cdot2}=\frac{\sqrt{3}}{3}\). 所以二面角 \(A_1-AB-B_1\) 的大小为 \(\arccos\frac{\sqrt{3}}{3}\).
\end{enumerate}
\end{solution}

\begin{problem}
已知正项数列 \(\{a_{n}\}\)，其前 \(n\) 项和 \(S_{n}\) 满足 \(10S_{n} = a_{n}^{2} + 5a_{n} + 6\)，且 \(a_{1}\)，\(a_{3}\)，\(a_{15}\) 成等比数列，求数列 \(\{a_{n}\}\) 的通项 \(a_{n}\).
\end{problem}

\begin{solution}
由 \(n=1\) 时的条件，得 \(10a_1=a_1^2+5a_1+6\)，即 \(a_1^2-5a_1+6=0\)，所以 \(a_1=2\) 或 \(a_1=3\). 当 \(n\geqslant2\) 时，将条件分别对 \(n\) 和 \(n-1\) 写出并相减，得 \(10a_n=a_n^2+5a_n-a_{n-1}^2-5a_{n-1}\)，即 \((a_n-a_{n-1})(a_n+a_{n-1}+5)=10a_n\). 因为数列各项为正数，所以右端为正，且 \(a_n+a_{n-1}+5>0\)，从而 \(a_n>a_{n-1}\). 进一步有 \(a_n^2-5a_n=a_{n-1}^2+5a_{n-1}\)，即 \(\left(a_n-\frac{5}{2}\right)^2=\left(a_{n-1}+\frac{5}{2}\right)^2\). 因此 \(a_n-\frac{5}{2}=\pm\left(a_{n-1}+\frac{5}{2}\right)\). 若取正号，则 \(a_n=a_{n-1}+5\)；若取负号，则 \(a_n=-a_{n-1}<0\)，与正项数列矛盾，所以对所有 \(n\geqslant2\)，都有 \(a_n=a_{n-1}+5\). 于是 \(a_n=a_1+5(n-1)\). 若 \(a_1=2\)，则 \(a_3=12\)，\(a_{15}=72\)，满足 \(a_3^2=12^2=2\cdot72=a_1a_{15}\)；若 \(a_1=3\)，则 \(a_3=13\)，\(a_{15}=73\)，而 \(a_3^2=169\neq219=a_1a_{15}\). 因为 \(a_1\)，\(a_3\)，\(a_{15}\) 为正项等比数列，必须满足 \(a_3^2=a_1a_{15}\)，故只能取 \(a_1=2\)，从而 \(a_n=5n-3\). 验证得 \(S_n=\sum_{k=1}^n(5k-3)=\frac{n(5n-1)}{2}\)，所以 \(10S_n=25n^2-5n=(5n-3)^2+5(5n-3)+6=a_n^2+5a_n+6\)，故所求通项为 \(a_n=5n-3\).
\end{solution}

\begin{problem}
如图，三定点 \(A(2,1)\)，\(B(0,-1)\)，\(C(-2,1)\)，三动点 \(D\)，\(E\)，\(M\) 满足 \(\overrightarrow{AD} = t\overrightarrow{AB}\)，\(\overrightarrow{BE} = t\overrightarrow{BC}\)，\(\overrightarrow{DM} = t\overrightarrow{DE}\)，\(t \in [0,1]\). 求：

\begin{enumerate}
\item 动直线 \(DE\) 斜率的变化范围；
\item 动点 \(M\) 的轨迹方程.
\end{enumerate}

\bitmapinlinefigure[9.6cm]{img/2006/shaanxi_science/q21_fig1.png}
\end{problem}

\begin{solution}
\begin{enumerate}
\item 由 \(\overrightarrow{AD}=t\overrightarrow{AB}\) 和 \(\overrightarrow{BE}=t\overrightarrow{BC}\)，得 \(D=(2-2t,1-2t)\)，\(E=(-2t,-1+2t)\). 因而 \(\overrightarrow{DE}=E-D=(-2,4t-2)\)，且 \(x_E-x_D=-2\neq0\)，所以动直线 \(DE\) 的斜率为 \(k=\frac{4t-2}{-2}=1-2t\). 当 \(t\in[0,1]\) 时，\(k\) 从 \(1\) 连续变化到 \(-1\)，故斜率的变化范围为 \([-1,1]\).
\item 由 \(\overrightarrow{DM}=t\overrightarrow{DE}\)，得 \(M=D+t(E-D)=(2-4t,1-4t+4t^2)\). 记 \(M\) 的坐标为 \((x,y)\)，则 \(x=2-4t=2(1-2t)\)，\(y=1-4t+4t^2=(1-2t)^2\)，所以 \(y=\frac{x^2}{4}\). 又因为 \(t\in[0,1]\)，故 \(-2\leqslant x\leqslant2\). 因此动点 \(M\) 的轨迹方程为 \(y=\frac{x^2}{4}\)，其中 \(-2\leqslant x\leqslant2\).
\end{enumerate}
\end{solution}

\begin{problem}
已知函数 \( f(x) = x^3 - x^2 + \frac{x}{2} + \frac{1}{4} \)，且存在 \( x_0 \in \left(0, \frac{1}{2}\right) \)，使 \( f(x_0) = x_0 \).

\begin{enumerate}
\item 证明： \( f(x) \) 是 \(\R\) 上的单调增函数；
\item 设 \(x_{1} = 0\)，\(x_{n + 1} = f(x_{n})\)，\(y_{1} = \frac{1}{2}\)，\(y_{n + 1} = f(y_{n})\)，其中 \(n = 1,2,\dots\)，证明：\(x_{n} < x_{n + 1} < x_0 < y_{n + 1} < y_{n}\)；
\item 证明： \( \frac{y_{n+1}-x_{n+1}}{y_{n}-x_{n}}<\frac{1}{2} \).
\end{enumerate}
\end{problem}

\begin{solution}
\begin{enumerate}
\item \(f'(x)=3x^2-2x+\frac{1}{2}=3\left(x-\frac{1}{3}\right)^2+\frac{1}{6}>0\)，所以 \(f(x)\) 在 \(\R\) 上单调递增.
\item 令 \(g(x)=f(x)-x=x^3-x^2-\frac{x}{2}+\frac{1}{4}\). 当 \(0\leqslant x\leqslant\frac{1}{2}\) 时，\(g'(x)=3x^2-2x-\frac{1}{2}\leqslant-\frac{x}{2}-\frac{1}{2}<0\)，所以 \(g(x)\) 在 \(\left[0,\frac{1}{2}\right]\) 上单调递减. 由 \(g(x_0)=0\)，可知当 \(0\leqslant x<x_0\) 时 \(g(x)>0\)，当 \(x_0<x\leqslant\frac{1}{2}\) 时 \(g(x)<0\). 初始时 \(x_1=0<x_0<\frac{1}{2}=y_1\). 若 \(x_n<x_0<y_n\)，则由 \(g(x_n)>0\)、\(f\) 的单调性和 \(g(y_n)<0\)，依次得到 \(x_n<f(x_n)<f(x_0)=x_0<f(y_n)<y_n\)，即 \(x_n<x_{n+1}<x_0<y_{n+1}<y_n\). 由数学归纳法，结论对所有 \(n=1,2,\dots\) 成立.
\item 由 \(x_{n+1}=f(x_n)\)、\(y_{n+1}=f(y_n)\)，有 \(\frac{y_{n+1}-x_{n+1}}{y_n-x_n}=\frac{f(y_n)-f(x_n)}{y_n-x_n}=x_n^2+x_ny_n+y_n^2-x_n-y_n+\frac{1}{2}\). 根据第（2）问，\(0\leqslant x_n<x_0<y_n\leqslant\frac{1}{2}\)，所以 \(0<x_n+y_n<1\)，且 \(x_n^2+x_ny_n+y_n^2=(x_n+y_n)^2-x_ny_n\leqslant(x_n+y_n)^2<x_n+y_n\). 因而 \(\frac{y_{n+1}-x_{n+1}}{y_n-x_n}<\frac{1}{2}\).
\end{enumerate}
\end{solution}
